\begin{table}[H]
\centering
\small
\caption{Monthly Macroeconomic Dynamics, Central Bank Solvency, and External Constraints across Structural Junctures (Chile, 1970--1973)}
\label{tab:marglin_monthly_grounding}
\label{tab:marglin_grounding}
\renewcommand{\arraystretch}{1.10}
\setlength{\tabcolsep}{4.8pt}
\resizebox{\textwidth}{!}{%
\begin{tabular}{lcccccc}
\toprule
& \textbf{Nov 1970} & \textbf{Aug 1971} & \textbf{Dec 1971} & \textbf{Oct 1972} & \textbf{Aug 1973} & \textbf{Dec 1973} \\
\textbf{Macroeconomic Variable} & \textit{Inauguration} & \textit{Gold Suspension} & \textit{Output Peak} & \textit{October Lockout} & \textit{Pre-Coup} & \textit{Post-Coup} \\
\midrule
\multicolumn{7}{l}{\textbf{A. Inflation and Price Dynamics}} \\[0.2em]
Monthly Inflation (\%) & 0.6\% & 1.1\% & 2.8\% & 15.2\% & 17.1\% & 4.7\% \\
YoY Inflation (\%) & 35.3\% & 17.4\% & 22.1\% & 142.9\% & 303.6\% & 508.1\% \\
\midrule
\multicolumn{7}{l}{\textbf{B. Real Physical Activity (Nov 1970 = 100)}} \\[0.2em]
Manufacturing Output Index & 100.0 & 122.0 & 133.3 & 111.9 & 104.6 & 118.1 \\
Mining Production Index & 100.0 & 104.9 & 107.0 & 114.7 & 100.7 & 119.6 \\
\midrule
\multicolumn{7}{l}{\textbf{C. Central Bank Solvency and Monetary Backing}} \\[0.2em]
Gross International Reserves (USD M) & 408.0 & 268.1 & 217.2 & 103.4 & 114.2 & 179.7 \\
Solvency Ratio Base ($\text{SolvR}^H \equiv \frac{e \cdot IR}{H \cdot 1000}$) & 0.63 & 0.22 & 0.17 & 0.07 & 0.06 & 0.25 \\
Solvency Ratio M1 ($\text{SolvR}^{M1} \equiv \frac{e \cdot IR}{M1 \cdot 1000}$) & 0.51 & 0.18 & 0.14 & 0.07 & 0.06 & 0.27 \\
Observed Multiplier ($M1/H$) & 1.23 & 1.23 & 1.23 & 1.03 & 0.98 & 0.95 \\
Base Money Monthly Growth ($g_H$) & 0.0\% & 11.5\% & 15.1\% & 12.6\% & 10.7\% & 21.8\% \\
M1 Money Monthly Growth ($g_{M1}$) & -1.6\% & 4.3\% & 10.2\% & 9.1\% & 11.6\% & 19.2\% \\
\midrule
\multicolumn{7}{l}{\textbf{D. Official Exchange Rate and World Commodity Prices}} \\[0.2em]
Official Exchange Rate ($e_t$) & 0.012 & 0.0122 & 0.0146 & 0.0250 & 0.0740 & 0.34 \\
World Copper Price (USD/lb) & \$0.492 & \$0.499 & \$0.471 & \$0.466 & \$0.948 & \$1.011 \\
World Gold Price (USD/oz) & \$37.44 & \$42.71 & \$43.48 & \$64.86 & \$106.76 & \$106.72 \\
\midrule
\multicolumn{7}{l}{\textbf{E. External Trade, Capacity to Import, and Structural Imbalance}} \\[0.2em]
Export Activity Index & 2.68 & 1.71 & 1.73 & 1.46 & 2.31 & 2.81 \\
Import Activity Index & 1.51 & 3.54 & 3.11 & 1.79 & 4.11 & 2.86 \\
Real Trade Ratio ($\text{Exp}/\text{Imp}$) & 1.78 & 0.48 & 0.56 & 0.82 & 0.56 & 0.98 \\
Real Capacity to Import ($M^{\text{cap}}_t$, Nov 1970 = 100) & 100.0 & 62.5 & 59.5 & 47.7 & 129.5 & 168.7 \\
Effective Real Imports ($M_t$, Nov 1970 = 100) & 100.0 & 234.5 & 206.2 & 118.4 & 272.3 & 189.4 \\
Real Structural Imbalance ($\Theta_t$, \%) & 100.0\% & 374.9\% & 346.4\% & 248.1\% & 210.3\% & 112.3\% \\
US Wholesale Price Index (Nov 1970 = 100) & 100.0 & 103.8 & 104.0 & 108.1 & 128.0 & 127.8 \\
\bottomrule
\end{tabular}%
}
\vspace{0.3em}
\begin{minipage}{\textwidth}
\footnotesize \textbf{Sources and Notes:} Primary monthly series from Central Bank of Chile and IMF International Financial Statistics. Gross international reserves ($IR_t$) represent official foreign exchange holdings from IMF IFS. Solvency Ratio defined relative to Base Money as $\text{SolvR}^{H}_t \equiv (e_t \cdot IR_t) / (H_t \cdot 1000)$ and relative to M1 as $\text{SolvR}^{M1}_t \equiv (e_t \cdot IR_t) / (M1_t \cdot 1000)$, where $IR_t$ is in USD millions, $e_t$ is the official exchange rate in pesos per USD (standardized per BCCh historical conventions), and $H_t, M1_t$ are recorded in thousands of millions of pesos (billions, \emph{miles de millones de pesos}), scaled by $1{,}000$ to express monetary liabilities in millions of pesos and ensure dimensional consistency with the numerator ($e_t \cdot IR_t$). Capacity to Import Index operationalizes the Prebisch-ECLAC formula: $M^{\text{cap}}_t \equiv (\text{Exports}_t / \text{Exports}_{70:11}) \times [(P^{\text{Cu}}_t / P^{\text{Cu}}_{70:11}) / (P^{\text{US}}_{W,t} / P^{\text{US}}_{W,70:11})] \times 100$. Real Structural Imbalance Ratio defined as $\Theta_t \equiv (M_t / M^{\text{cap}}_t) \times 100$, tracking the physical quantum of imports relative to real export purchasing power. Selected dates mark historical milestones: Nov 1970 (Allende inauguration baseline); Aug 1971 (Eximbank credit denial and suspension of dollar-gold convertibility); Dec 1971 (cyclical output peak); Oct 1972 (national trucking strike and commercial lockout); Aug 1973 (late administration baseline before the September coup); and Dec 1973 (post-coup benchmark following October 1973 Decree Law 522 price deregulation).
\end{minipage}
\end{table}
